\documentclass[11pt,a4paper]{article}
\usepackage[margin=1in]{geometry}
\usepackage{booktabs}
\usepackage{amsmath}
\usepackage{graphicx}
\usepackage{hyperref}
\usepackage{xcolor}
\usepackage{float}
\usepackage[font=small]{caption}
\usepackage{enumitem}
\usepackage{parskip}

\title{German National PV Generation Forecasting\\
       \large Technical Note --- Internship Assessment}
\author{Xiyun Fu}
\date{\today}

\begin{document}
\maketitle

%% =====================================================================
\section{Data Sources}
%% =====================================================================

Table~\ref{tab:sources} summarises all data sources considered.
The pipeline uses three operational sources---Open-Meteo (ERA5 reanalysis and ICON-D2 NWP forecasts), SMARD actual generation, and BNetzA PV capacity---plus pvlib for deterministic solar features.

\begin{table}[H]
\centering
\caption{Data sources considered and used.}\label{tab:sources}
\small
\begin{tabular}{@{}llllp{4.2cm}@{}}
\toprule
Source & Variables & Resolution & Status & Rationale \\
\midrule
Open-Meteo ERA5 & GHI, $T_{2m}$, $v_{10}$, RH & Hourly, 0.25\textdegree & \textbf{Used} & Free, no auth, covers 2024--2025 \\
Open-Meteo ICON-D2 & GHI, $T_{2m}$, cloud & 15-min, ${\sim}$2\,km & \textbf{Used} & Native 15-min NWP for target-time features \\
SMARD (BNetzA) & Actual solar gen. & 15-min & \textbf{Used} & Free CSV/JSON, no API key \\
BNetzA / MaStR & PV capacity & Static (2024) & \textbf{Used} & Aggregated per Bundesland \\
pvlib & Clear-sky, zenith & Computed & \textbf{Used} & Deterministic solar geometry \\
\midrule
ERA5 via CDS API & ssrd, t2m, tcc & Hourly, 0.25\textdegree & Rejected & CDS registration delay; Open-Meteo wraps ERA5 identically \\
ENTSO-E Transparency & Actual PV gen. & 15-min & Rejected & Requires API key; SMARD provides identical data freely \\
MaStR bulk export & Plant-level PV & Static & Rejected & 500\,MB download; aggregated BNetzA totals suffice \\
OSM Overpass & PV locations & Static & Rejected & Incomplete coverage; capacity weights from BNetzA preferred \\
\bottomrule
\end{tabular}
\end{table}

\subsection{Data Overview}

\textbf{Weather (ERA5).}
280,704 rows (16 states $\times$ 17,544 hours) covering 2024-01-01 to 2025-12-31.
Variables: \texttt{shortwave\_radiation} (GHI, W/m$^2$), \texttt{temperature\_2m} (\textdegree C), \texttt{wind\_speed\_10m} (m/s), \texttt{relative\_humidity\_2m} (\%).
Served via \texttt{archive-api.open-meteo.com/v1/archive}.

\textbf{NWP Forecasts (ICON-D2).}
1,122,816 rows (16 states $\times$ 70,176 quarter-hours).
Native 15-min resolution at ${\sim}$2\,km.
Variables: GHI, temperature, cloud cover.
Served via \texttt{historical-forecast-api.open-meteo.com/v1/forecast} with \texttt{models=icon\_d2}.
ICON-D2 updates every 3\,h with ${\sim}$1.5--2\,h publication delay, so forecasts for time $T{+}h$ are available well before prediction time $T$.

\textbf{Target (SMARD).}
70,176 rows of quarter-hourly German national solar generation (MW), 2024-01-01 to 2025-12-31.
Range: 0.5--13,209\,MW; daytime mean: 2,106\,MW.

\textbf{PV Capacity.}
Static 2024 snapshot from BNetzA: 80,400\,MWp total across 16 Bundesl\"ander.
Bayern dominates (27.1\%), followed by Baden-W\"urttemberg (12.7\%) and NRW (11.1\%).

%% =====================================================================
\section{Cleaning and Harmonisation}
%% =====================================================================

The harmonisation pipeline (\texttt{harmonise.py}) produces a unified 15-min modelling table:

\begin{enumerate}[nosep]
  \item \textbf{Weather interpolation.} Hourly ERA5 $\to$ 15-min via cubic spline for GHI (non-negative constraint) and linear interpolation for temperature, wind speed, and humidity. Per-state processing preserves spatial structure.
  \item \textbf{Capacity-weighted national averages.} For variable $x$:
  \begin{equation}
    x^{\text{nat}}_t = \sum_{i=1}^{16} w_i \cdot x_{i,t}, \qquad w_i = \frac{C_i}{\sum_j C_j}
  \end{equation}
  where $C_i$ is installed MWp in state $i$.  Applied to both ERA5 and ICON-D2 data.
  \item \textbf{NWP merge.} ICON-D2 data (already at 15-min) is joined on timestamp via left merge.  If the NWP parquet is absent, this step is skipped gracefully.
  \item \textbf{Target join.} Inner join with SMARD actual generation on timestamp.
  \item \textbf{Gap filling.} Forward-fill up to 4\,h (16 intervals); remaining NaN rows dropped.
  \item \textbf{Type casting.} All numeric columns cast to float32 (${\sim}$7.6\,MB on disk for 70,173 rows $\times$ 26 columns).
\end{enumerate}

\begin{table}[H]
\centering
\caption{Modelling table summary statistics (70,173 rows, zero nulls).}\label{tab:stats}
\small
\begin{tabular}{@{}lrrr@{}}
\toprule
Column & Min & Max & Mean \\
\midrule
\texttt{ghi\_wm2\_national} (W/m$^2$) & 0.0 & 889.4 & 138.7 \\
\texttt{temperature\_2m\_national} (\textdegree C) & $-$7.5 & 34.8 & 10.7 \\
\texttt{wind\_speed\_10m\_national} (m/s) & 2.0 & 33.4 & 10.8 \\
\texttt{nwp\_ghi\_wm2\_national} (W/m$^2$) & 0.0 & 883.3 & 127.6 \\
\texttt{nwp\_cloud\_cover\_national} (\%) & 0.0 & 100.0 & 73.7 \\
\texttt{actual\_solar\_mw} (MW) & 0.5 & 13,209 & 1,951 \\
\texttt{proxy\_solar\_mw} (MW) & 0.0 & 12,022 & 1,875 \\
\bottomrule
\end{tabular}
\end{table}

%% =====================================================================
\section{Synthetic Proxy Construction}
%% =====================================================================

The synthetic proxy (\texttt{build\_proxy.py}) aggregates state-level irradiance into a single national generation estimate:

\begin{equation}
P^{\text{raw}}_t = \sum_{i=1}^{16} w_i \cdot \text{GHI}_{i,t}, \qquad
P^{\text{syn}}_t = \eta \cdot P^{\text{raw}}_t
\end{equation}

where $w_i$ is the normalised capacity weight for state $i$ and $\eta$ is a scaling factor fitted via OLS (no intercept) on daytime training data (2024, where $P^{\text{raw}} > 0$ and actual $> 0$).

\textbf{Proxy quality} on the full dataset (daytime hours):
correlation $r = 0.968$, $R^2 = 0.937$, MAE $= 490$\,MW, nMAE $= 0.214$.
The proxy serves as a floor reference; both learned models must beat it.

%% =====================================================================
\section{Modelling}
%% =====================================================================

\subsection{Forecasting Paradigm}

At each quarter-hour $T$, the model predicts generation for $T{+}1$ through $T{+}16$ (15\,min to 4\,h ahead).  The feature matrix is split into:

\begin{table}[H]
\centering
\caption{Feature availability at prediction time $T$.}\label{tab:features}
\small
\begin{tabular}{@{}lll@{}}
\toprule
Type & Available for & Examples \\
\midrule
Origin (time $T$) & Past only & ERA5 GHI/temp/wind, proxy, lags, clearsky index \\
Target deterministic ($T{+}h$) & All times & Clear-sky GHI, solar zenith, hour/DoY sin/cos \\
Target NWP ($T{+}h$) & Forecast & ICON-D2 GHI, temperature, cloud cover \\
\bottomrule
\end{tabular}
\end{table}

Origin features come from the prediction moment; target features come from the predicted moment.
Both are causally valid: deterministic features are computable for any time; NWP forecasts were issued hours before $T$.
Lag features use $\geq$24\,h shifts (\texttt{actual\_lag\_1d} = shift 96, \texttt{actual\_lag\_7d} = shift 672) to respect ENTSO-E's ${\sim}$60-min publication delay.
A night mask zeroes all predictions where target solar zenith $> 85$\textdegree.

\subsection{Baseline Model (Ridge Regression)}

Ridge regression ($\alpha = 1.0$) with \texttt{StandardScaler} on 6 features:
\texttt{proxy\_solar\_mw}, \texttt{clearsky\_index}, \texttt{forecast\_horizon},
\texttt{target\_clearsky\_ghi}, \texttt{target\_solar\_zenith}, \texttt{target\_nwp\_ghi\_wm2\_national}.

\subsection{Candidate Model (LightGBM)}

Gradient-boosted trees with 18 features (7 origin + 7 target deterministic + 3 target NWP + horizon).

\textbf{Hyperparameter tuning.}
Optuna (30 trials, 5-min timeout) with 3-fold rolling-origin CV within 2024:
Fold 1: train Jan--Jun, val Jul--Aug;
Fold 2: train Jan--Aug, val Sep--Oct;
Fold 3: train Jan--Oct, val Nov--Dec.

Tuned hyperparameters: \texttt{num\_leaves}$\,{=}\,$53, \texttt{learning\_rate}$\,{=}\,$0.0142, \texttt{min\_child\_samples}$\,{=}\,$32, \texttt{subsample}$\,{=}\,$0.532, \texttt{colsample\_bytree}$\,{=}\,$0.729, \texttt{reg\_alpha}$\,{=}\,$1.57, \texttt{reg\_lambda}$\,{=}\,$8.99, \texttt{n\_estimators}$\,{=}\,$800 with early stopping (50 rounds).

\subsection{Walk-Forward Evaluation}

Twelve monthly rounds simulate operational deployment over 2025:

\begin{quote}
\textit{Round $r$:} train on Jan~2024 through month $r{-}1$ of 2025; validate on the last month of the training window (early stopping); test on month $r$ of 2025.
\end{quote}

Each round retrains from scratch using the Optuna-tuned HPs.
A ``frozen'' model (Round~1, trained only on 2024) is also evaluated on all months to quantify the benefit of retraining.

%% =====================================================================
\section{Validation}
%% =====================================================================

\subsection{Temporal Split}

\begin{table}[H]
\centering
\caption{Data split (walk-forward scheme).}\label{tab:split}
\small
\begin{tabular}{@{}llr@{}}
\toprule
Split & Period & Rows (multi-horizon) \\
\midrule
Training (Round 1) & Jan 2024 -- Dec 2024 & ${\sim}$560{,}000 \\
Test (all 12 rounds) & Jan 2025 -- Dec 2025 & 560{,}456 \\
\bottomrule
\end{tabular}
\end{table}

\subsection{Leakage Controls}

\begin{enumerate}[nosep]
  \item \textbf{No future weather.} Origin features (ERA5 GHI, temperature, wind) are taken from time $T$ only.  For $T{+}h$, only deterministic and NWP features are used.
  \item \textbf{No future actuals.} Lag features start at $T{-}96$ (24\,h ago) and $T{-}672$ (7\,d ago). No \texttt{actual\_lag\_1} through \texttt{actual\_lag\_3} exist.
  \item \textbf{Scaler fitted on train only.} \texttt{StandardScaler.fit\_transform()} on training data; \texttt{.transform()} on validation/test.
  \item \textbf{No HP tuning on test.} Optuna uses 2024 data only; 2025 is never seen during tuning.
  \item \textbf{Strict temporal ordering.} Walk-forward ensures train cutoff $<$ test start in every round.
\end{enumerate}

\subsection{Metrics}

\begin{equation}
\text{MAE} = \frac{1}{N}\sum_{t=1}^{N} |\hat{y}_t - y_t|, \qquad
\text{nMAE} = \frac{\text{MAE}}{\bar{y}_{\text{daytime}}}, \qquad
\text{Skill} = 1 - \frac{\text{MAE}_{\text{model}}}{\text{MAE}_{\text{baseline}}}
\end{equation}

where $\bar{y}_{\text{daytime}}$ is the mean actual generation during hours with $y_t > 0$ (2,106\,MW).

%% =====================================================================
\section{Results}
%% =====================================================================

\subsection{Overall Performance}

\begin{table}[H]
\centering
\caption{Overall test metrics (all horizons, 560{,}456 predictions).}\label{tab:overall}
\small
\begin{tabular}{@{}lrrrr@{}}
\toprule
Model & MAE (MW) & RMSE (MW) & nMAE & Skill vs BL \\
\midrule
Synthetic proxy & 490 & --- & 0.214 & --- \\
Baseline (Ridge) & 298 & 589 & 0.141 & --- \\
LightGBM (retrained) & 254 & 494 & 0.121 & 0.146 \\
LightGBM (frozen) & 265 & 510 & 0.126 & 0.111 \\
\bottomrule
\end{tabular}
\end{table}

The retrained LightGBM reduces nMAE by 14.6\% relative to the baseline.
Monthly retraining provides a further 4.0\% improvement over the frozen model, confirming the value of expanding the training window.

\subsection{Monthly Walk-Forward Results}

\begin{table}[H]
\centering
\caption{Per-round nMAE (walk-forward, all horizons).}\label{tab:monthly}
\small
\begin{tabular}{@{}clrrr@{}}
\toprule
Round & Month & Baseline & LightGBM & Frozen \\
\midrule
1 & Jan & 0.226 & 0.290 & 0.290 \\
2 & Feb & 0.222 & 0.177 & 0.180 \\
3 & Mar & 0.227 & 0.167 & 0.169 \\
4 & Apr & 0.127 & 0.125 & 0.127 \\
5 & May & 0.089 & 0.098 & 0.100 \\
6 & Jun & 0.087 & 0.084 & 0.089 \\
7 & Jul & 0.115 & 0.087 & 0.086 \\
8 & Aug & 0.109 & 0.088 & 0.090 \\
9 & Sep & 0.182 & 0.116 & 0.117 \\
10 & Oct & 0.194 & 0.146 & 0.152 \\
11 & Nov & 0.221 & 0.196 & 0.232 \\
12 & Dec & 0.263 & 0.259 & 0.349 \\
\bottomrule
\end{tabular}
\end{table}

LightGBM underperforms the baseline in Round~1 (January) because the model has no 2025 winter data yet; the baseline's simpler structure generalises better from the 2024 training set.
From Round~2 onward, LightGBM consistently outperforms.
The frozen model degrades sharply in Nov--Dec (nMAE 0.232/0.349 vs retrained 0.196/0.259), demonstrating that retraining captures seasonal adaptation.

\subsection{Horizon Analysis}

Error increases modestly with forecast lead time, from nMAE~0.120 at $h{=}1$ (15\,min) to 0.122 at $h{=}16$ (4\,h).
The small gradient ($+$1.7\% over 4 hours) reflects that NWP forecast features for the target time partially compensate for staleness of origin features.

\begin{table}[H]
\centering
\caption{LightGBM nMAE by forecast horizon (selected steps).}\label{tab:horizon}
\small
\begin{tabular}{@{}rrr@{}}
\toprule
Horizon & Minutes & nMAE \\
\midrule
1 & 15 & 0.1195 \\
4 & 60 & 0.1202 \\
8 & 120 & 0.1206 \\
12 & 180 & 0.1211 \\
16 & 240 & 0.1215 \\
\bottomrule
\end{tabular}
\end{table}

\subsection{Feature Importance}

\begin{table}[H]
\centering
\caption{Top 10 features by LightGBM gain (single-split model).}\label{tab:importance}
\small
\begin{tabular}{@{}clr@{}}
\toprule
Rank & Feature & Gain ($\times 10^{12}$) \\
\midrule
1 & \texttt{proxy\_solar\_mw} & 30,266 \\
2 & \texttt{clearsky\_ghi} & 1,564 \\
3 & \texttt{actual\_lag\_1d} & 1,199 \\
4 & \texttt{hour\_sin} & 682 \\
5 & \texttt{doy\_sin} & 183 \\
6 & \texttt{clearsky\_index} & 173 \\
7 & \texttt{solar\_zenith} & 152 \\
8 & \texttt{doy\_cos} & 127 \\
9 & \texttt{wind\_speed\_10m\_national} & 109 \\
10 & \texttt{temperature\_2m\_national} & 52 \\
\bottomrule
\end{tabular}
\end{table}

The synthetic proxy dominates ($>$85\% of total gain), confirming that capacity-weighted irradiance is the primary driver.
Clear-sky GHI and the 24-hour lag provide complementary persistence and astronomical signals.

%% =====================================================================
\section{Limitations and Next Steps}
%% =====================================================================

\textbf{Winter underperformance.}
nMAE exceeds 0.25 in Nov--Dec (Table~\ref{tab:monthly}), driven by low solar elevation and high cloud variability.
Satellite-derived irradiance (e.g., EUMETSAT MTG, 10-min latency) could improve near-real-time cloud state estimation in winter.

\textbf{Round~1 cold start.}
LightGBM's nMAE of 0.290 in January 2025 is worse than the baseline (0.226), indicating the model overfits to 2024 patterns.
A regularisation schedule or domain-adaptive pretraining could mitigate this.

\textbf{Flat horizon curve.}
The 1.7\% nMAE increase from 15\,min to 4\,h suggests origin features (observed weather, lags) contribute less than expected.
For longer horizons ($>$4\,h), this signal will degrade further.
Recursive or autoregressive strategies with NWP-only features should be explored for day-ahead forecasting ($h{=}96$).

\textbf{ERA5 is not real-time.}
ERA5 reanalysis has a ${\sim}$5-day publication lag.
In operations, observed weather at time $T$ would come from station networks, satellite products, or short-range NWP nowcasts---not ERA5.
The pipeline already integrates ICON-D2 NWP; extending to live feeds is straightforward.

\textbf{Deterministic forecasts only.}
The pipeline produces point estimates.
Quantile regression (LightGBM with \texttt{objective=quantile}) or conformal prediction intervals would better serve trading decisions under uncertainty.

\textbf{Single national aggregate.}
Regional disaggregation (e.g., per-TSO-zone) with high-resolution ICON-D2 inputs would enable locational marginal pricing and grid congestion applications.

\end{document}
